\chapter{Corporate Bonds}\label{ch:m2:corporate-bonds}

A company's bonds are rated BBB$-$, the lowest rung of \omterm{def:m2:corporate-bonds:rating}{investment grade}. For
months its bonds' spread over Treasuries has been widening, as analysts and
investors read the same numbers the rating agencies read. Then one agency cuts
it to BB$+$. The bonds are due to leave the investment-grade indices; funds whose
mandates allow only \omterm{def:m2:corporate-bonds:rating}{investment grade} must sell them, and high-yield funds, which did not own them yesterday, must decide at what price
to buy. When the European Central Bank studied such ``\omterm{def:m2:corporate-bonds:rating}{fallen angels}'', it found
that most of the damage to their prices came before the downgrade, and that
after the last agency had downgraded them their prices partly recovered, as if
the forced sales had pushed them too far. This chapter introduces corporate
bonds: how they are issued, what their ratings mean, the several spreads by
which their prices are quoted and compared, and the \omterm{def:m2:corporate-bonds:covenant}{covenants} and \omterm{def:m2:corporate-bonds:covenant}{seniority} that
decide what their holders recover when a company fails.

\section{Issuance and the new-issue concession}

\begin{definition}[New-issue concession]\label{def:m2:corporate-bonds:nic}
The \emph{new-issue concession}\index{new-issue concession} is the extra spread
at which a company's new bond is priced over its existing bonds of similar
maturity, the price the issuer pays to attract enough buyers quickly.
\end{definition}

A company issues bonds through a syndicate of banks that underwrite them (One
Quant Book 1, chapter 2): they announce the deal, take orders from investors over a
few hours at an initial price guidance, tighten the price as the book of orders
grows, and allocate. The bonds are priced as a spread over a benchmark Treasury of
similar maturity. When a new issue comes wide of where the company's outstanding
bonds trade, the difference is the concession: a reward for absorbing new supply
and for committing large amounts at once. A concession that disappears in the first
days of trading is the investors' gain and the issuer's cost.

\section{Ratings and the investment-grade line}

\begin{definition}[Credit rating, investment grade, high yield, fallen angel]\label{def:m2:corporate-bonds:rating}
A \emph{credit rating}\index{credit rating} is a rating agency's opinion of an
issuer's or a bond's creditworthiness, on a letter scale from AAA (Aaa) down to D.
Bonds rated BBB$-$ (Baa3) or above are \emph{investment grade}\index{investment grade};
those rated BB$+$ (Ba1) or below are \emph{high yield}\index{high yield}. A
\emph{fallen angel}\index{fallen angel} is an issuer downgraded from investment
grade to high yield by at least one of the three major agencies.
\end{definition}

The line between BBB$-$ and BB$+$ is one notch on a scale, but it divides the
market. Many insurers, pension funds and bond funds may hold only \omterm{def:m2:corporate-bonds:rating}{investment
grade}; indices are split at the line, and so are the funds that track them;
regulators and central banks use it too. When the Federal Reserve set up its
facility to buy newly issued corporate bonds in 2020, it required issuers to be
rated at least BBB$-$/Baa3 on 22 March 2020, and allowed those downgraded afterwards
down to BB$-$/Ba3: the line, with a bridge for the fallen.

The ECB's study found the market's pattern around the line: credit default swap
premiums of \omterm{def:m2:corporate-bonds:rating}{fallen angels} widened well before the first downgrade, moved little
on the day, and partly recovered after the last investment-grade rating was lost;
funds tracking indices mostly use sampling and may keep a downgraded bond for a
while, which spreads the selling over time. Of the \omterm{def:m2:corporate-bonds:rating}{fallen angels} in Moody's data
over twenty years, nearly a quarter returned to \omterm{def:m2:corporate-bonds:rating}{investment grade}, almost half
stayed in \omterm{def:m2:corporate-bonds:rating}{high yield} and 12\% defaulted. \Cref{fig:m2:corporate-bonds:angel}
draws the pattern on an illustrative bond.

\begin{omfigure}
\begin{tikzpicture}
\begin{axis}[width=11cm,height=4.8cm,
  xlabel={months from the downgrade that takes the bond out of investment grade},ylabel={spread (\%)},
  tick label style={font=\small},label style={font=\small},
  xmin=-6.5,xmax=6.5,ymin=2.3,ymax=4.3,grid=major,grid style={gray!25},
  table/col sep=comma]
\addplot[omThm,very thick,mark=*,mark size=1.6pt] table[x=month,y=spread]{figdata/markets-2/21-corporate-bonds/angel.csv};
\draw[omDef,dashed] (axis cs:-6.5,3.5) -- (axis cs:6.5,3.5);
\node[omDef,font=\scriptsize,anchor=south west] at (axis cs:-6.3,3.52) {fair for BB$+$};
\draw[gray,dashed] (axis cs:0,2.3) -- (axis cs:0,4.3);
\node[omProp,font=\scriptsize,anchor=west] at (axis cs:0.15,4.08) {forced selling};
\end{axis}
\end{tikzpicture}

\omcaption{The spread of an illustrative \omterm{def:m2:corporate-bonds:rating}{fallen angel}: most of the widening comes
before the downgrade, a forced-selling overshoot on the day, then a recovery towards
the level fair for its new rating, the pattern the ECB found in the credit default
swaps of \omterm{def:m2:corporate-bonds:rating}{fallen angels}. Illustrative path; data: the chapter's
tutorial.}\label{fig:m2:corporate-bonds:angel}
\end{omfigure}

\section{The family of spread measures}

\begin{definition}[Credit spread; G-, I-, Z- and asset-swap spreads]\label{def:m2:corporate-bonds:spreads}
A \emph{credit spread}\index{credit spread} is the extra yield a risky bond pays over
a reference curve. The \emph{G-spread}\index{G-spread} is the bond's yield minus the
government yield interpolated at its maturity; the \emph{I-spread}\index{I-spread}
its yield minus the interpolated swap rate. The \emph{Z-spread}\index{Z-spread} is
the constant spread which, added to the swap zero curve, discounts the bond's cash
flows to its price. The \emph{asset-swap spread}\index{asset-swap spread} is the
spread over the floating rate that a buyer of the bond earns by swapping its fixed
coupons into floating payments, at par: the difference between the bond's value
on the swap curve and its price, divided by the annuity of the floating leg.
\end{definition}

The four measures answer different questions. The \omterm{def:m2:corporate-bonds:spreads}{G-spread} compares the bond with
the safest asset; it mixes credit with the \omterm{def:m2:interest-rate-swaps:spread}{swap spread} and with any mismatch of
coupon and curve shape. The \omterm{def:m2:corporate-bonds:spreads}{I-spread} compares it with swaps, the instrument dealers
hedge with. The \omterm{def:m2:corporate-bonds:spreads}{Z-spread} uses the whole zero curve, not one point, and so prices
the bond's actual cash flows. The \omterm{def:m2:corporate-bonds:spreads}{asset-swap spread} is what an investor who funds
at the floating rate and buys the bond with a swap actually earns. On a flat curve
and a bond at par they coincide; in general they do not.

\begin{proposition}[Why the measures differ]\label{prop:m2:corporate-bonds:differ}
For a bond whose yield is $y$ and whose maturity is $T$: the G- and \omterm{def:m2:corporate-bonds:spreads}{I-spreads} differ
by the \omterm{def:m2:interest-rate-swaps:spread}{swap spread} at $T$; the \omterm{def:m2:corporate-bonds:spreads}{Z-spread} differs from the \omterm{def:m2:corporate-bonds:spreads}{I-spread} by the effect of
the curve's slope on the bond's cash flows, which a single yield ignores; and the
\omterm{def:m2:corporate-bonds:spreads}{asset-swap spread} and the \omterm{def:m2:corporate-bonds:spreads}{Z-spread} measure the same gap between the bond's price
and its value on the swap curve, the first as a running spread paid on par, the
second as a shift of every discount rate: they agree to first order near par, and
part as the spread grows and the price moves away from par.
\end{proposition}

\begin{proof}
The first follows from the definitions. A yield discounts every cash flow at one
rate; the \omterm{def:m2:corporate-bonds:spreads}{Z-spread} discounts each at the zero rate of its date plus a constant;
when the curve slopes, the two agree only if the zero rates equal the yield.
The asset swap turns the gap $P_{\mathrm{curve}} - P$ into a running spread by
dividing it by the floating leg's annuity $A$. A small \omterm{def:m2:corporate-bonds:spreads}{Z-spread} $z$ lowers the price
by about $z\sum_i t_i c_i P(0,t_i)$, a duration-weighted sum close to $100A$ for a
bond near par; so the two spreads are close. The \omterm{def:m2:corporate-bonds:spreads}{Z-spread}'s effect on price is not
linear, and the weights differ when the coupon is far from the swap rate, so the gap
widens with the spread and with the distance from par.
\end{proof}

\begin{example}[One bond, four spreads]\label{ex:m2:corporate-bonds:four}
A seven-year corporate bond pays a 5.50\% semiannual coupon and trades at 99.00:
its yield is 5.6751\%. With the illustrative curves of
\cref{fig:m2:corporate-bonds:curves}, the seven-year Treasury at 4.50\% and swap at
4.25\%, its \omterm{def:m2:corporate-bonds:spreads}{G-spread} is 117.5 basis points and its \omterm{def:m2:corporate-bonds:spreads}{I-spread} 142.5. Its \omterm{def:m2:corporate-bonds:spreads}{Z-spread} over
the swap zero curve is 140.7 basis points and its \omterm{def:m2:corporate-bonds:spreads}{asset-swap spread} 142.9.
\end{example}

\begin{omfigure}
\begin{tikzpicture}
\begin{axis}[width=11cm,height=5cm,
  xlabel={maturity (years)},ylabel={yield (\%)},
  tick label style={font=\small},label style={font=\small},
  xmin=0,xmax=31,ymin=3.9,ymax=5.9,grid=major,grid style={gray!25},
  legend columns=3,legend style={font=\footnotesize,at={(0.5,-0.3)},anchor=north,draw=none,fill=none,/tikz/every even column/.append style={column sep=8pt}},
  table/col sep=comma]
\addplot[omDef,very thick] table[x=t,y=govt]{figdata/markets-2/21-corporate-bonds/curves.csv};
\addplot[omThm,very thick,dashed] table[x=t,y=swap]{figdata/markets-2/21-corporate-bonds/curves.csv};
\addplot[omProp,only marks,mark=*,mark size=2.4pt] table[x=t,y=y]{figdata/markets-2/21-corporate-bonds/bond.csv};
\draw[omProp,thick,{Stealth}-{Stealth}] (axis cs:7.6,4.50) -- (axis cs:7.6,5.675);
\node[omProp,font=\scriptsize,anchor=west] at (axis cs:7.9,5.05) {G: 117.5 bp};
\draw[omThm,thick,{Stealth}-{Stealth}] (axis cs:6.4,4.25) -- (axis cs:6.4,5.675);
\node[omThm,font=\scriptsize,anchor=east] at (axis cs:6.1,4.85) {I: 142.5 bp};
\legend{Treasury par curve,swap par curve,the bond}
\end{axis}
\end{tikzpicture}

\omcaption{The bond of \cref{ex:m2:corporate-bonds:four} against illustrative
Treasury and swap curves. The \omterm{def:m2:corporate-bonds:spreads}{G-spread} is measured to the Treasury curve and the
\omterm{def:m2:corporate-bonds:spreads}{I-spread} to the swap curve, which lies below it here, as it has at long maturities
since the \omterm{def:m2:interest-rate-swaps:spread}{swap spreads} turned negative (\cref{ch:m2:interest-rate-swaps}).
Illustrative; data: the chapter's tutorial.}\label{fig:m2:corporate-bonds:curves}
\end{omfigure}

\begin{definition}[Option-adjusted spread]\label{def:m2:corporate-bonds:oas}
The \emph{option-adjusted spread}\index{option-adjusted spread} (OAS) of a bond with
embedded options is the spread over the curve at which the bond's value, computed
with a model of rates that values the options, equals its price: the \omterm{def:m2:corporate-bonds:spreads}{Z-spread} less
the options' cost expressed as a spread.
\end{definition}

A \omterm{def:m2:the-rates-options-market:swaption}{callable bond}'s \omterm{def:m2:corporate-bonds:spreads}{Z-spread} overstates what its holder is paid for credit, because
part of the yield is payment for the call the holder has sold to the issuer
(\cref{ch:m2:the-rates-options-market}). If the bond of
\cref{ex:m2:corporate-bonds:four} were callable at par in three years, the call
would be a \omterm{def:m2:the-rates-options-market:swaption}{receiver swaption} on the remaining four years struck at the coupon.
Valued at 100 basis points of normal volatility it is worth 4.80 points, 79.7
basis points a year over the bond's life: an \omterm{def:m2:corporate-bonds:oas}{option-adjusted spread} of about 61
basis points against a \omterm{def:m2:corporate-bonds:spreads}{Z-spread} of 141.

The history of \omterm{def:m2:corporate-bonds:spreads}{credit spreads} is a history of crises. The Treasury publishes a
high-quality corporate yield curve, built from bonds rated AAA, AA and A; its
ten-year rate over the ten-year Treasury yield rose from under 1 percentage point in
early 2007 to 5.04 in October 2008, and to about 2.1 in the spring of 2020
(\cref{fig:m2:corporate-bonds:hqm}).

\begin{omfigure}
\begin{tikzpicture}
\begin{axis}[width=11cm,height=4.8cm,
  ylabel={percentage points},
  tick label style={font=\small},label style={font=\small},
  xmin=1984,xmax=2026.75,ymin=0,ymax=5.5,grid=major,grid style={gray!25},
  xtick={1985,1990,1995,2000,2005,2010,2015,2020,2025},xticklabel style={/pgf/number format/1000 sep={}},
  table/col sep=comma]
\addplot[omThm,thick] table[x=t,y=spread]{figdata/markets-2/21-corporate-bonds/hqm.csv};
\end{axis}
\end{tikzpicture}

\omcaption{Ten-year high-quality corporate spot rate (Treasury HQM curve, bonds
rated AAA, AA or A) minus the ten-year Treasury yield, monthly, 1984 to August 2026.
The measure mixes credit with the difference between a zero-coupon and a par rate,
but its peaks, 2008 and 2020, are the credit crises. Data: FRED series HQMCB10YR
(US Treasury) and GS10.}\label{fig:m2:corporate-bonds:hqm}
\end{omfigure}

\begin{dated}{2026-09}{Spreads now}\label{dat:m2:corporate-bonds:spreads}
Ten-year HQM corporate spot rate 5.58\% and ten-year Treasury 4.68\% in August 2026
(monthly averages): a gap of 0.90 percentage points. The widely quoted index spreads
of corporate bonds are published by index providers under licence and are not
reproduced here.
\end{dated}

\section{Covenants, seniority and recovery}

\begin{definition}[Covenant, seniority, recovery rate]\label{def:m2:corporate-bonds:covenant}
A \emph{covenant}\index{covenant} is a promise in a bond's terms that restricts what
the issuer may do (borrow more, pay dividends, sell assets, merge) or requires it
to maintain some condition. \emph{Seniority}\index{seniority} is a claim's rank in
the order of payment in insolvency: secured before senior unsecured before
subordinated before equity. The \emph{recovery rate}\index{recovery rate} is the share
of a bond's face value its holders receive after the issuer defaults.
\end{definition}

\begin{omfigure}
\centering
\begin{tikzpicture}[font=\small,x=1cm,y=0.42cm]
\fill[omDef!25] (0,0) rectangle (2.4,6); \node[font=\footnotesize,align=center] at (1.2,3) {firm value\\at default\\600};
\fill[omThm!30] (4,0) rectangle (6.4,2); \node[font=\footnotesize] at (5.2,1) {secured 200};
\fill[omProp!30] (4,2) rectangle (6.4,7); \node[font=\footnotesize,align=center] at (5.2,4.5) {senior\\unsecured 500};
\fill[gray!25] (4,7) rectangle (6.4,9); \node[font=\footnotesize] at (5.2,8) {subordinated 200};
\draw[thick,dashed] (2.6,6) -- (7.2,6);
\node[font=\footnotesize,anchor=west,align=left] at (7.3,1) {paid in full: 100\%};
\node[font=\footnotesize,anchor=west,align=left] at (7.3,4.2) {share 400 of 500: 80\%};
\node[font=\footnotesize,anchor=west,align=left] at (7.3,8) {nothing left: 0\%};
\end{tikzpicture}

\omcaption{\omterm{def:m2:corporate-bonds:covenant}{Seniority} and recovery on an illustrative default. A firm worth 600 at
default owes 200 secured, 500 senior unsecured and 200 subordinated: the secured
creditors are paid in full, the senior unsecured share the remaining 400, and the
subordinated recover nothing. Schematic.}\label{fig:m2:corporate-bonds:seniority}
\end{omfigure}

Investment-grade bonds usually carry few \omterm{def:m2:corporate-bonds:covenant}{covenants}; high-yield bonds carry more,
because their holders need protection against a company that takes on more debt
or pays its cash to shareholders. \omterm{def:m2:corporate-bonds:covenant}{Seniority} decides recovery: in a liquidation or
restructuring, senior secured claims are paid first out of their collateral, senior
unsecured claims share what is left, and subordinated claims are paid only after
them (\cref{fig:m2:corporate-bonds:seniority}). \omterm{def:m2:corporate-bonds:covenant}{Recovery rates} vary widely with the firm, the industry and the cycle, which
is why a spread is compensation for two uncertain numbers: the probability of
default and the loss when it happens, the subject of
\cref{ch:m2:credit-default-swaps}.

\begin{remark}[Spread, default and loss]\label{rem:m2:corporate-bonds:loss}
Over a year, a spread $s$ roughly pays for an annual default probability $p$ times a
loss given default $1 - R$: $s \approx p(1 - R)$, plus premia for risk and
illiquidity. A spread of 140 basis points with a recovery of 40\% corresponds, if it
were all expected loss, to a default probability of 2.3\% a year. Whatever the
bonds' actual default rate falls short of that is the premium investors demand for
bearing default risk and illiquidity.
\end{remark}

\section{Tutorial: four spreads of one bond}\label{tut:m2:corporate-bonds:spreads}

\begin{tutorial}
\textbf{Goal.} Compute the G-, I-, Z- and \omterm{def:m2:corporate-bonds:spreads}{asset-swap spreads} of a bond, and its
\omterm{def:m2:corporate-bonds:oas}{option-adjusted spread} if it is callable; plot the spread history.
\textbf{End state:}
\cref{fig:m2:corporate-bonds:curves,fig:m2:corporate-bonds:hqm},
\cref{ex:m2:corporate-bonds:four} and the numbers of the weekend problem.
\begin{tutsteps}
\item \textbf{Curves and yields}: interpolation, the bond's cash flows, price and
yield.
\omcode{code/firm/spreads/firm_spreads.py}{10}{33}{Curve interpolation, cash flows, price and yield.}\label{lst:m2:corporate-bonds:yield}
\item \textbf{The spreads}: to par curves, to the zero curve, and the asset swap.
\omcode{code/firm/spreads/firm_spreads.py}{36}{61}{G-, I-, Z- and asset-swap spreads.}\label{lst:m2:corporate-bonds:spreads}
\item \textbf{Run} \texttt{credit\_demo.measures()},
\texttt{credit\_demo.callable\_oas()}, \texttt{credit\_demo.fallen\_angel()} and
\texttt{fig\_credit.py}.
\end{tutsteps}
\textbf{What to change next.} Steepen the swap curve and watch the \omterm{def:m2:corporate-bonds:spreads}{Z-spread} move
away from the \omterm{def:m2:corporate-bonds:spreads}{I-spread}; then price the bond at 110 and at 90 and compare the
asset-swap and \omterm{def:m2:corporate-bonds:spreads}{Z-spreads}, as \cref{prop:m2:corporate-bonds:differ} predicts.
\end{tutorial}

\section{Build: the spread-measure library}\label{bld:m2:corporate-bonds:spreads}

\begin{build}
\bfield{Purpose} The miniature firm quotes and compares corporate bonds, hedges
them with Treasuries or swaps, and must speak each client's spread language.
\bfield{Interface} \texttt{interp(curve, t)}; \texttt{flows}; \texttt{price\_from\_yield};
\texttt{yield\_from\_price}; \texttt{g\_spread}; \texttt{i\_spread};
\texttt{price\_on\_zero}; \texttt{z\_spread}; \texttt{asset\_swap\_spread};
\texttt{spread\_price\_impact(duration, widening, price)}.
\bfield{Rules} Curves as (years, rate) points, linearly interpolated; semiannual
coupons; prices dirty on a coupon date; zero rates continuously compounded; root
finding by bisection.
\bfield{Acceptance tests} \texttt{code/firm/spreads/tests/}: interpolation and yield
round trips; spreads equal on flat curves; the \omterm{def:m2:corporate-bonds:spreads}{Z-spread} recovered from a price built
with it; an \omterm{def:m2:corporate-bonds:spreads}{asset-swap spread} of zero for a bond priced on the curve.
\bfield{Stretch} Accrued interest and settlement dates (the bond build of
\cref{ch:m2:government-bonds}); OAS for Bermudan calls with a short-rate lattice
(One Quant Book 6); spreads to a fitted issuer curve.
\end{build}

\begin{omsources}
\item M.~Belloni, T.~Helmersson, M.~Jarmuzek, B.~Mosk and F.~Nikolic,
``Understanding what happens when angels fall'', ECB \emph{Financial Stability Review},
November 2020.
\item Federal Reserve, Primary Market Corporate Credit Facility term sheet, April 2020.
\item US Department of the Treasury, High Quality Market corporate bond yield curve;
FRED series HQMCB10YR and GS10.
\end{omsources}

\section{Exercises}

\begin{exercise}[$\star$]\label{exo:m2:corporate-bonds:1}
A ten-year bond yields 6.10\%. The ten-year Treasury yields 4.68\% and the ten-year
swap rate is 4.40\%. Give its G- and \omterm{def:m2:corporate-bonds:spreads}{I-spreads}.
\end{exercise}

\begin{exercise}[$\star$]\label{exo:m2:corporate-bonds:2}
Rated Baa3 by one agency and BB$+$ by another, is a bond \omterm{def:m2:corporate-bonds:rating}{investment grade}?
Is its issuer a \omterm{def:m2:corporate-bonds:rating}{fallen angel} if it was BBB$-$ at both last month?
\end{exercise}

\begin{exercise}[$\star$]\label{exo:m2:corporate-bonds:3}
A bond of \omterm{def:m2:government-bonds:duration}{modified duration} 6 sees its spread widen by 80 basis points. Give its
price change per 100.
\end{exercise}

\begin{exercise}[$\star\star$]\label{exo:m2:corporate-bonds:4}
In \cref{ex:m2:corporate-bonds:four}, the \omterm{def:m2:corporate-bonds:spreads}{Z-spread} and the \omterm{def:m2:corporate-bonds:spreads}{asset-swap spread} differ
by 2.2 basis points. Why are they close, and when would they differ by more?
\end{exercise}

\begin{exercise}[$\star\star$]\label{exo:m2:corporate-bonds:5}
A spread of 140 basis points and a recovery of 40\%: what annual default probability
would the spread pay for if it were all expected loss? What else does the spread pay
for?
\end{exercise}

\begin{exercise}[$\star\star$]\label{exo:m2:corporate-bonds:6}
Why does a new issue come with a concession, and who gains if it disappears?
\end{exercise}

\begin{exercise}[$\star\star\star$]\label{exo:m2:corporate-bonds:7}
\emph{Coding.} With the build, reprice the bond of \cref{ex:m2:corporate-bonds:four}
at 105 and give its four spreads.
\end{exercise}

\begin{exercise}[$\star\star\star$]\label{exo:m2:corporate-bonds:8}
\emph{Find the flaw.} ``This \omterm{def:m2:the-rates-options-market:swaption}{callable bond} yields 140 basis points over swaps, the
same as the bullet bond of the same issuer, so it is equally cheap.'' Correct it.
\end{exercise}

\section{Problem: The Fallen Angel}

\begin{problem}[{Weekend problem --- what forced selling costs and pays}]\label{pb:m2:corporate-bonds:1}
A fund holds USD~50 million of a company's bonds, \omterm{def:m2:government-bonds:duration}{modified duration} 6. Six months
before any downgrade their spread is 250 basis points; it widens to 290 three months
before, 330 one month before, and 400 on the day the last agency cuts the company to
BB$+$ and the bonds leave the investment-grade indices; then 380 a month later, 350
after three months and 345 after six. Analysts judge 350 fair for a BB$+$ credit.

\medskip\noindent\textbf{Part I --- Before the downgrade.}
\begin{enumerate}
\item What did the bonds lose, per 100, between six months before and the downgrade?
\item How much of that happened before the downgrade day?
\item Why does the market move before the agencies?
\item What would an investment-grade fund have done with a negative outlook?
\item Give the reaction on the day, per 100.
\end{enumerate}

\medskip\noindent\textbf{Part II --- The forced sale.}
\begin{enumerate}[resume]
\item Give the overshoot over fair value, in spread and per 100.
\item What does the overshoot cost the fund if it sells on the day?
\item Who buys, and why do they need a discount?
\item How do index rules and fund sampling change the pressure?
\item What does the fund gain by waiting three months, if it may?
\end{enumerate}

\medskip\noindent\textbf{Part III --- The buyer.}
\begin{enumerate}[resume]
\item Give a buyer's gain per 100 from buying on the day and selling three months later.
\item What risk does the buyer bear?
\item What do the ECB's figures on later outcomes say about that risk?
\item How does the carry of the higher spread add to the return?
\item Why is the trade not free money?
\end{enumerate}

\medskip\noindent\textbf{Part IV --- Judgement.}
\begin{enumerate}[resume]
\item Why do ratings matter so much when prices move first?
\item How would a central-bank facility like that of 2020 change the pattern?
\item What would you ask the fund's mandate to allow?
\item State the \emph{named result}: the forced-selling discount and its recovery.
\item In one sentence: why do \omterm{def:m2:corporate-bonds:rating}{fallen angels} overshoot?
\end{enumerate}
\end{problem}

\section{Interview questions}

\begin{interviewq}[$\star$ \iqroles{trader, researcher}]\label{iq:m2:corporate-bonds:1}
What is the difference between the \omterm{def:m2:corporate-bonds:spreads}{G-spread}, the \omterm{def:m2:corporate-bonds:spreads}{Z-spread} and the \omterm{def:m2:corporate-bonds:spreads}{asset-swap spread}?
\end{interviewq}

\begin{interviewq}[$\star$ \iqroles{trader, bank}]\label{iq:m2:corporate-bonds:2}
What is a \omterm{def:m2:corporate-bonds:rating}{fallen angel}, and why do its bonds often cheapen around the downgrade?
\end{interviewq}

\begin{interviewq}[$\star\star$ \iqroles{researcher}]\label{iq:m2:corporate-bonds:3}
How would you decompose a \omterm{def:m2:corporate-bonds:spreads}{credit spread} into expected loss and risk premium?
\end{interviewq}

\begin{interviewq}[$\star\star$ \iqroles{trader}]\label{iq:m2:corporate-bonds:4}
A client asks for a price on USD~50 million of a BBB bond that rarely trades. How do
you price it?
\end{interviewq}

\begin{interviewq}[$\star\star$ \iqroles{researcher, trader}]\label{iq:m2:corporate-bonds:5}
Why does the OAS of a \omterm{def:m2:the-rates-options-market:swaption}{callable bond} differ from its \omterm{def:m2:corporate-bonds:spreads}{Z-spread}, and which should you
compare across bonds?
\end{interviewq}

\begin{interviewq}[$\star\star\star$ \iqroles{developer, researcher}]\label{iq:m2:corporate-bonds:6}
Design a system that computes spreads for fifty thousand corporate bonds every minute.
\end{interviewq}
